%% ═══════════════════════════════════════════════════════════════════════════
%% INTRODUCTION — included via %% ═══════════════════════════════════════════════════════════════════════════
%% INTRODUCTION — included via %% ═══════════════════════════════════════════════════════════════════════════
%% INTRODUCTION — included via %% ═══════════════════════════════════════════════════════════════════════════
%% INTRODUCTION — included via \input{latex/introduction} from main.tex
%% ═══════════════════════════════════════════════════════════════════════════

\section{Introduction}\label{sec:introduction}

Cross-sectional momentum, the tendency for recent stock market winners to continue outperforming and losers to continue underperforming, is one of the most robust anomalies in empirical finance \citep{JegadeeshTitman1993}. Yet it is also one of the most fragile: momentum portfolios suffered losses exceeding 50\% during the 2009 market rebound \citep{DanielMoskowitz2016}. This tension between long-run profitability and vulnerability to regime shifts motivates the central question of this thesis: how does a real-time regime signal interact with the cross-sectional momentum anomaly, and what does this interaction reveal about momentum's underlying structure?

Two prominent responses address momentum's fragility. \citet{BarrosoSantaClara2015} scale exposure inversely with recent volatility; \citet{GHM2023} blend momentum horizons according to a market-state classification derived from trailing returns. Both improve upon unconditional momentum, but both rely on backward-looking indicators to infer the current regime. This thesis proposes an alternative: a Bayesian Hidden Markov Model estimated on real-time stress indicators (drawdown, dispersion, relative market participation, and credit spread) produces a filtered probability of market panic, which is fed alongside 12 momentum horizons into machine learning models that learn how to adapt cross-sectional stock selection to the current regime.

The approach differs from prior work in three respects. First, the regime classification uses real-time stress indicators rather than trailing returns. Second, the regime signal enters the model as one feature among many, allowing the data to determine its optimal use rather than imposing a fixed rule. Third, by comparing a linear model (logistic regression) with a nonlinear model (XGBoost) on identical features, the framework provides a clean test of whether the regime-momentum interaction requires nonlinear modelling. The central claim is that the regime signal is not a direct return predictor but a \emph{context variable} -- a feature that governs which momentum horizons are informative in a given market state.

Out-of-sample evaluation on 167 months of U.S. equity data (2011--2025) yields three main findings. First, only XGBoost produces meaningful long-short returns (Sharpe 1.08, six-factor alpha 24.7\%, $t = 4.78$); logistic regression with identical features collapses (Sharpe $-$0.03), establishing that the interaction is inherently nonlinear. Second, the regime signal is essential -- removing it halves performance -- and most valuable during stress, with the long-short spread widening from 1.3\% to 3.1\% per month in panic. Third, model-free information coefficient analysis confirms this rotation in the raw data: all 12 momentum horizons are significant in calm, none in panic, and the short horizons flip sign.

Specifically, this thesis addresses three questions: (i)~does the regime signal's interaction with momentum require nonlinear modelling, and if so, why does a linear model fail? (ii)~how does the cross-sectional structure of momentum change between calm and panic regimes? (iii)~how robust is the regime-momentum interaction to alternative signal specifications, risk preferences, and prolonged market stress?

The remainder of the thesis is organised as follows. Section~\ref{sec:literature} reviews the momentum anomaly, regime-switching models, market-state momentum, and machine learning in asset pricing. Section~\ref{sec:methodology} develops the two-stage framework: the Bayesian HMM for regime classification and the three cross-sectional stock selection methods. Section~\ref{sec:data} describes the data, HMM feature selection procedure, and cross-sectional features. Section~\ref{sec:results} presents the out-of-sample results, including regime-conditional performance, stress tests, SHAP analysis, ablation studies, and the risk-aversion analysis. Section~\ref{sec:conclusion} summarises the findings, discusses limitations, and suggests directions for future research.
 from main.tex
%% ═══════════════════════════════════════════════════════════════════════════

\section{Introduction}\label{sec:introduction}

Cross-sectional momentum, the tendency for recent stock market winners to continue outperforming and losers to continue underperforming, is one of the most robust anomalies in empirical finance \citep{JegadeeshTitman1993}. Yet it is also one of the most fragile: momentum portfolios suffered losses exceeding 50\% during the 2009 market rebound \citep{DanielMoskowitz2016}. This tension between long-run profitability and vulnerability to regime shifts motivates the central question of this thesis: how does a real-time regime signal interact with the cross-sectional momentum anomaly, and what does this interaction reveal about momentum's underlying structure?

Two prominent responses address momentum's fragility. \citet{BarrosoSantaClara2015} scale exposure inversely with recent volatility; \citet{GHM2023} blend momentum horizons according to a market-state classification derived from trailing returns. Both improve upon unconditional momentum, but both rely on backward-looking indicators to infer the current regime. This thesis proposes an alternative: a Bayesian Hidden Markov Model estimated on real-time stress indicators (drawdown, dispersion, relative market participation, and credit spread) produces a filtered probability of market panic, which is fed alongside 12 momentum horizons into machine learning models that learn how to adapt cross-sectional stock selection to the current regime.

The approach differs from prior work in three respects. First, the regime classification uses real-time stress indicators rather than trailing returns. Second, the regime signal enters the model as one feature among many, allowing the data to determine its optimal use rather than imposing a fixed rule. Third, by comparing a linear model (logistic regression) with a nonlinear model (XGBoost) on identical features, the framework provides a clean test of whether the regime-momentum interaction requires nonlinear modelling. The central claim is that the regime signal is not a direct return predictor but a \emph{context variable} -- a feature that governs which momentum horizons are informative in a given market state.

Out-of-sample evaluation on 167 months of U.S. equity data (2011--2025) yields three main findings. First, only XGBoost produces meaningful long-short returns (Sharpe 1.08, six-factor alpha 24.7\%, $t = 4.78$); logistic regression with identical features collapses (Sharpe $-$0.03), establishing that the interaction is inherently nonlinear. Second, the regime signal is essential -- removing it halves performance -- and most valuable during stress, with the long-short spread widening from 1.3\% to 3.1\% per month in panic. Third, model-free information coefficient analysis confirms this rotation in the raw data: all 12 momentum horizons are significant in calm, none in panic, and the short horizons flip sign.

Specifically, this thesis addresses three questions: (i)~does the regime signal's interaction with momentum require nonlinear modelling, and if so, why does a linear model fail? (ii)~how does the cross-sectional structure of momentum change between calm and panic regimes? (iii)~how robust is the regime-momentum interaction to alternative signal specifications, risk preferences, and prolonged market stress?

The remainder of the thesis is organised as follows. Section~\ref{sec:literature} reviews the momentum anomaly, regime-switching models, market-state momentum, and machine learning in asset pricing. Section~\ref{sec:methodology} develops the two-stage framework: the Bayesian HMM for regime classification and the three cross-sectional stock selection methods. Section~\ref{sec:data} describes the data, HMM feature selection procedure, and cross-sectional features. Section~\ref{sec:results} presents the out-of-sample results, including regime-conditional performance, stress tests, SHAP analysis, ablation studies, and the risk-aversion analysis. Section~\ref{sec:conclusion} summarises the findings, discusses limitations, and suggests directions for future research.
 from main.tex
%% ═══════════════════════════════════════════════════════════════════════════

\section{Introduction}\label{sec:introduction}

Cross-sectional momentum, the tendency for recent stock market winners to continue outperforming and losers to continue underperforming, is one of the most robust anomalies in empirical finance \citep{JegadeeshTitman1993}. Yet it is also one of the most fragile: momentum portfolios suffered losses exceeding 50\% during the 2009 market rebound \citep{DanielMoskowitz2016}. This tension between long-run profitability and vulnerability to regime shifts motivates the central question of this thesis: how does a real-time regime signal interact with the cross-sectional momentum anomaly, and what does this interaction reveal about momentum's underlying structure?

Two prominent responses address momentum's fragility. \citet{BarrosoSantaClara2015} scale exposure inversely with recent volatility; \citet{GHM2023} blend momentum horizons according to a market-state classification derived from trailing returns. Both improve upon unconditional momentum, but both rely on backward-looking indicators to infer the current regime. This thesis proposes an alternative: a Bayesian Hidden Markov Model estimated on real-time stress indicators (drawdown, dispersion, relative market participation, and credit spread) produces a filtered probability of market panic, which is fed alongside 12 momentum horizons into machine learning models that learn how to adapt cross-sectional stock selection to the current regime.

The approach differs from prior work in three respects. First, the regime classification uses real-time stress indicators rather than trailing returns. Second, the regime signal enters the model as one feature among many, allowing the data to determine its optimal use rather than imposing a fixed rule. Third, by comparing a linear model (logistic regression) with a nonlinear model (XGBoost) on identical features, the framework provides a clean test of whether the regime-momentum interaction requires nonlinear modelling. The central claim is that the regime signal is not a direct return predictor but a \emph{context variable} -- a feature that governs which momentum horizons are informative in a given market state.

Out-of-sample evaluation on 167 months of U.S. equity data (2011--2025) yields three main findings. First, only XGBoost produces meaningful long-short returns (Sharpe 1.08, six-factor alpha 24.7\%, $t = 4.78$); logistic regression with identical features collapses (Sharpe $-$0.03), establishing that the interaction is inherently nonlinear. Second, the regime signal is essential -- removing it halves performance -- and most valuable during stress, with the long-short spread widening from 1.3\% to 3.1\% per month in panic. Third, model-free information coefficient analysis confirms this rotation in the raw data: all 12 momentum horizons are significant in calm, none in panic, and the short horizons flip sign.

Specifically, this thesis addresses three questions: (i)~does the regime signal's interaction with momentum require nonlinear modelling, and if so, why does a linear model fail? (ii)~how does the cross-sectional structure of momentum change between calm and panic regimes? (iii)~how robust is the regime-momentum interaction to alternative signal specifications, risk preferences, and prolonged market stress?

The remainder of the thesis is organised as follows. Section~\ref{sec:literature} reviews the momentum anomaly, regime-switching models, market-state momentum, and machine learning in asset pricing. Section~\ref{sec:methodology} develops the two-stage framework: the Bayesian HMM for regime classification and the three cross-sectional stock selection methods. Section~\ref{sec:data} describes the data, HMM feature selection procedure, and cross-sectional features. Section~\ref{sec:results} presents the out-of-sample results, including regime-conditional performance, stress tests, SHAP analysis, ablation studies, and the risk-aversion analysis. Section~\ref{sec:conclusion} summarises the findings, discusses limitations, and suggests directions for future research.
 from main.tex
%% ═══════════════════════════════════════════════════════════════════════════

\section{Introduction}\label{sec:introduction}

Cross-sectional momentum, the tendency for recent stock market winners to continue outperforming and losers to continue underperforming, is one of the most robust anomalies in empirical finance \citep{JegadeeshTitman1993}. Yet it is also one of the most fragile: momentum portfolios suffered losses exceeding 50\% during the 2009 market rebound \citep{DanielMoskowitz2016}. This tension between long-run profitability and vulnerability to regime shifts motivates the central question of this thesis: how does a real-time regime signal interact with the cross-sectional momentum anomaly, and what does this interaction reveal about momentum's underlying structure?

Two prominent responses address momentum's fragility. \citet{BarrosoSantaClara2015} scale exposure inversely with recent volatility; \citet{GHM2023} blend momentum horizons according to a market-state classification derived from trailing returns. Both improve upon unconditional momentum, but both rely on backward-looking indicators to infer the current regime. This thesis proposes an alternative: a Bayesian Hidden Markov Model estimated on real-time stress indicators (drawdown, dispersion, relative market participation, and credit spread) produces a filtered probability of market panic, which is fed alongside 12 momentum horizons into machine learning models that learn how to adapt cross-sectional stock selection to the current regime.

The approach differs from prior work in three respects. First, the regime classification uses real-time stress indicators rather than trailing returns. Second, the regime signal enters the model as one feature among many, allowing the data to determine its optimal use rather than imposing a fixed rule. Third, by comparing a linear model (logistic regression) with a nonlinear model (XGBoost) on identical features, the framework provides a clean test of whether the regime-momentum interaction requires nonlinear modelling. The central claim is that the regime signal is not a direct return predictor but a \emph{context variable} -- a feature that governs which momentum horizons are informative in a given market state.

Out-of-sample evaluation on 167 months of U.S. equity data (2011--2025) yields three main findings. First, only XGBoost produces meaningful long-short returns (Sharpe 1.08, six-factor alpha 24.7\%, $t = 4.78$); logistic regression with identical features collapses (Sharpe $-$0.03), establishing that the interaction is inherently nonlinear. Second, the regime signal is essential -- removing it halves performance -- and most valuable during stress, with the long-short spread widening from 1.3\% to 3.1\% per month in panic. Third, model-free information coefficient analysis confirms this rotation in the raw data: all 12 momentum horizons are significant in calm, none in panic, and the short horizons flip sign.

Specifically, this thesis addresses three questions: (i)~does the regime signal's interaction with momentum require nonlinear modelling, and if so, why does a linear model fail? (ii)~how does the cross-sectional structure of momentum change between calm and panic regimes? (iii)~how robust is the regime-momentum interaction to alternative signal specifications, risk preferences, and prolonged market stress?

The remainder of the thesis is organised as follows. Section~\ref{sec:literature} reviews the momentum anomaly, regime-switching models, market-state momentum, and machine learning in asset pricing. Section~\ref{sec:methodology} develops the two-stage framework: the Bayesian HMM for regime classification and the three cross-sectional stock selection methods. Section~\ref{sec:data} describes the data, HMM feature selection procedure, and cross-sectional features. Section~\ref{sec:results} presents the out-of-sample results, including regime-conditional performance, stress tests, SHAP analysis, ablation studies, and the risk-aversion analysis. Section~\ref{sec:conclusion} summarises the findings, discusses limitations, and suggests directions for future research.
